\oldpage[II]{51}
$X_f\cap\varphi^{-1}(U)=D_+(f|U)$ dans $\varphi^{-1}(U)$ identifié à
$X_U=\operatorname{Proj}(\Gamma(U,\mathcal{S}))$ ($\varphi$ désignant le
morphisme structural $X\to Y$). En outre, le préschéma induit sur $X_f$ est
affine au-dessus de $Y$ et est canoniquement isomorphe à
$\operatorname{Spec}(\mathcal{S}^{(d)}/(f-1)\mathcal{S}^{(d)})$
\textup{(\hyperref[II.1.3.1-fr]{1.3.1})}.
\end{proposition}

\begin{proof}
On a $f|U\in\Gamma(U,\mathcal{S}_d)=(\Gamma(U,\mathcal{S}))_d$. Si $U$,
$U'$ sont deux ouverts affines de $Y$ tels que $U'\subset U$, $f|U'$ est
l'image de $f|U$ par l'homomorphisme de restriction
\[
  \Gamma(U,\mathcal{S})\longrightarrow\Gamma(U',\mathcal{S}),
\]
donc $D_+(f|U')$ est égal (avec les notations de
(\hyperref[II.3.1.1-fr]{3.1.1})) au préschéma induit sur l'image réciproque
$\rho_{U',U}^{-1}(D_+(f|U))$ dans $X_{U'}$
(\hyperref[II.2.8.1-fr]{2.8.1})~; d'où la première assertion. En outre, le
préschéma induit sur $D_+(f|U)$ par $X_U$ s'identifie canoniquement à
$\operatorname{Spec}((\Gamma(U,\mathcal{S}))_{f|U})$, ces identifications
étant compatibles avec les homomorphismes de restriction
(\hyperref[II.2.8.1-fr]{2.8.1})~; la seconde assertion résulte alors de
(\hyperref[II.2.2.5-fr]{2.2.5}) et de la commutativité du diagramme
(\hyperref[II.2.8.2.1-fr]{2.8.2.1}).
\end{proof}

On dira encore que $X_f$ (en tant qu'ouvert dans l'espace sous-jacent $X$)
est l'ensemble des $x\in X$ où $f$ \emph{ne s'annule pas}.

\begin{corollary}[3.1.5]
\label{II.3.1.5-fr}
Si $f\in\Gamma(Y,\mathcal{S}_d)$, $g\in\Gamma(Y,\mathcal{S}_e)$, on a
\[
\label{II.3.1.5.1-fr}
  X_{fg}=X_f\cap X_g.
\tag{3.1.5.1}
\]
\end{corollary}

\begin{proof}
Il suffit en effet de considérer l'intersection des deux membres avec un
ensemble $\varphi^{-1}(U)$, où $U$ est un ouvert affine dans $Y$, et
d'appliquer la formule (\hyperref[II.2.3.3.2-fr]{2.3.3.2}).
\end{proof}

\begin{corollary}[3.1.6]
\label{II.3.1.6-fr}
Soit $(f_\alpha)$ une famille de sections de $\mathcal{S}$ au-dessus de $Y$
telle que $f_\alpha\in\Gamma(Y,\mathcal{S}_{d_\alpha})$~; si le faisceau
d'idéaux de $\mathcal{S}$ engendré par cette famille
\textup{(\hyperref[0.5.1.1-fr]{\textbf{0}, 5.1.1})} contient tous les $\mathcal{S}_n$ à partir
d'un certain rang, l'espace sous-jacent $X$ est réunion des $X_{f_\alpha}$.
\end{corollary}

\begin{proof}
En effet, pour tout ouvert affine $U$ de $Y$, $\varphi^{-1}(U)$ est réunion
des $X_{f_\alpha}\cap\varphi^{-1}(U)$
(\hyperref[II.2.3.14-fr]{2.3.14}).
\end{proof}

\begin{corollary}[3.1.7]
\label{II.3.1.7-fr}
Soit $\mathcal{A}$ une $\mathcal{O}_Y$-Algèbre quasi-cohérente~; posons
\[
  \mathcal{S}=\mathcal{A}[T]
  =\mathcal{A}\otimes_{\mathbf{Z}}\mathbf{Z}[T]
\]
où $T$ est une indéterminée ($\mathbf{Z}$ et $\mathbf{Z}[T]$ étant
considérés comme faisceaux simples sur $Y$). Alors
$X=\operatorname{Proj}(\mathcal{S})$ s'identifie canoniquement à
$\operatorname{Spec}(\mathcal{A})$. En particulier,
$\operatorname{Proj}(\mathcal{O}_Y[T])$ s'identifie à $Y$.
\end{corollary}

\begin{proof}
En appliquant (\hyperref[II.3.1.6-fr]{3.1.6}) à l'unique section
$f\in\Gamma(Y,\mathcal{S})$ égale à $T$ en chaque point de $Y$, on voit que
$X_f=X$. En outre, on a ici $d=1$, et
$\mathcal{S}^{(1)}/(f-1)\mathcal{S}^{(1)}
=\mathcal{S}/(f-1)\mathcal{S}$ est canoniquement isomorphe à $\mathcal{A}$,
d'où le corollaire (\hyperref[II.1.2.2-fr]{1.2.2}).
\end{proof}

Soit $g\in\Gamma(Y,\mathcal{O}_Y)$~; si on prend
$\mathcal{S}=\mathcal{O}_Y[T]$, on a donc $g\in\Gamma(Y,\mathcal{S}_0)$~; soit
\[
  h=gT\in\Gamma(Y,\mathcal{S}_1).
\]
Si $X=\operatorname{Proj}(\mathcal{S})$, l'identification canonique définie
dans (\hyperref[II.3.1.7-fr]{3.1.7}) identifie $X_h$ à l'ensemble ouvert
$Y_g$ de $Y$ (au sens de (\hyperref[0.5.5.2-fr]{\textbf{0}, 5.5.2}))~: en effet, on
peut se borner au cas où $Y=\operatorname{Spec}(A)$ est affine, et tout se
ramène alors (compte tenu de (\hyperref[II.2.2.5-fr]{2.2.5})) au fait que
l'anneau de fractions $A_g$ s'identifie canoniquement à
$A[T]/(gT-1)A[T]$ (\hyperref[0.1.2.3-fr]{\textbf{0}, 1.2.3}).

\begin{proposition}[3.1.8]
\label{II.3.1.8-fr}
Soit $\mathcal{S}$ une $\mathcal{O}_Y$-Algèbre graduée quasi-cohérente à
degrés positifs.
\begin{enumerate}
  \item[\rm(i)] Pour tout $d>0$, il existe un $Y$-isomorphisme canonique de
    $\operatorname{Proj}(\mathcal{S})$ sur
    $\operatorname{Proj}(\mathcal{S}^{(d)})$.

